\subsection{Distance Sampling}

Distance sampling is a statistical framework by which the local population density and abundance
of a species is estimated.
Two primary methods, 'point' and 'line' transect distance sampling,\cite{distance_sampling_applications}
exist for this; however, the methodology explored in this project leverages the latter approach.

Here, it is assumed that the true distribution of individuals in the population is evenly distributed
over the survey area.
In principle, animal detections are surveyed along the length of a line (transect line) in the
direction of observation, originating at the observation point (origin).
For all observed animals, the animal's distance with respect to the origin is recorded.
An assumption is made that the probability of detecting an animal at distance zero with respect to
the origin is 1 and that this probability of detection decreases as the animal's distance from the
origin increases.
Therefore, to accurately derive an estimate of population density, the detectability of the population
with respect to distance must be modelled to account for undetected animals.

\subsubsection{Density and Abundance Derivation}

The following summarises the key steps for deriving density and abundance estimates according
to the comprehensive distance sampling methodology introduced in 'Distance Sampling: Estimating Abundance
of Biological Populations' (Buckland et al., 1993).\cite{distance_samping}

\vspace{5mm}
\textbf{The Detection Function}

The detection function, $g(y)$, describes the probability that a given animal will be detected
at distance $y$ from the origin.
The observed distances are grouped into distance interval bins that cover the range of observed distances
(e.g., 0–1 m, 1–2 m, \ldots) up to the furthest considered distance, $w$, known as the truncation distance.
To estimate density, a detection function is fitted to model the observed distribution of detection
distances over the transect line.

[Figure]

The 'wellness of fit' of several detection function candidates is assessed.
A variety of candidates can be used; however, common candidates include the half-normal and
hazard-rate distributions, given by
\[
    g(y) = e^{-\frac{y^2}{2\sigma^2}} \qquad \mathrm{and} \qquad g(y) = 1 - e^{-(\frac{y}{\sigma})^{-\beta}}
\]
respectively, where $\sigma$ and $\beta$ are determined parameters that maximise the wellness of fit
of the given candidate function to the distribution of observed distances.
To determine the parameterisation of a candidate function, these parameters are varied and the wellness
of fit of the specific parameterisation is assessed.

\vspace{5mm}
\textbf{The Probability Density Function}

To do this, the probability density function, $f(y)$, is derived for the parameterised candidate function,
given by
\[
    f(y) = \dfrac{g(y)}{\mu} \mathrm{,} \qquad \mathrm{where} \qquad \mu = \int_{0}^{w} g(y) \, dy \mathrm{.}
\]
The probability density function is a measure of the distribution of observed detections, i.e.,
given an animal is detected between zero and $w$ meters, the probability of the detection distance
being $y$ meters.

\vspace{5mm}
\textbf{Wellness of Fit Assessment}

The likelihood, $L$, of the probability density function, $f(y)$, is then calculated to determine the wellness
of fit of the corresponding detection function, $g(y)$.
This is given by
\[
    L = \prod_{i=1}^n f(y_i)\mathrm{,}
\]
where $y_i$ is the $i^{\mathrm{th}}$ observed distance and $n$ is the total number of observations.
Here, the likelihood gives a measure of the likelihood of the observed distribution of distances
arising, given that true detectability is modelled by the fitted detection function, $g(y)$.
As such, the appropriate parameters of the candidate function are determined to be those that result
in the maximisation of the corresponding likelihood, $L$.
Parameters for each of the candidate functions are determined using this approach.

Now with a set of parameterised candidate functions, the final detection function, $g(y)$, can be
determined by calculating the Akaike information criterion, AIC, for each of the candidate functions.
This is given by
\[
    \mathrm{AIC} = 2q - 2ln(L) \mathrm{,}
\]
where $L$ is the likelihood resulting from the given candidate function and $q$ is the number of
parameters of the given candidate function.
Here, $ln(L)$, is referred to as the 'log-likelihood' and gives a measure of the wellness of fit of
the given candidate function.
Conversely, $q$ is a measure of complexity of the given candidate function and contributes to a
'penalty' term in the equation.
A such, the final detection function, $g(y)$, is determined to be the candidate function that gives
the lowest AIC value.

\vspace{5mm}
\textbf{Density and Abundance Estimates}

Density, $D$, is then given by
\[
    D = \dfrac{n}{2 \mu L} \mathrm{,} \qquad \mathrm{where} \qquad \mu = \int_{0}^{w} g(y) \, dy \mathrm{,}
\]
where $n$ is the total number of observed individuals, $L$ is the total length of the transect line
and $w$ is the truncation distance.
Abundance, $N$, is then given by
\[
    N = D \cdot A \mathrm{,}
\]
where $D$ is density and $A$ is area.



\subsubsection{Software Implementation}

